\documentclass[12pt, a4paper]{article}

% ============================================================
%  Packages
% ============================================================
\usepackage[utf8]{inputenc}
\usepackage[T1]{fontenc}
\usepackage[english]{babel}
\usepackage{amsmath, amssymb}
\usepackage[ruled, vlined, linesnumbered]{algorithm2e}
\usepackage{geometry}
\usepackage{hyperref}

\geometry{margin=2.5cm}

% ============================================================
%  Document metadata
% ============================================================
\title{Discrete Cuckoo Search for the CARP --- Pseudocode}
\author{}
\date{}

% ============================================================
\begin{document}
\maketitle

% ------------------------------------------------------------
\section{Introduction and biological analogy}
% ------------------------------------------------------------

The \textbf{Cuckoo Search} algorithm (Yang \& Deb, 2009) is inspired by the
brood parasitism of the cuckoo bird: the cuckoo lays its eggs in the nests of
other birds and, if the egg is not detected as an intruder, the chick is raised
by the host. In optimization terms:

\begin{itemize}
    \item Each \textbf{nest} is a candidate solution (a set of CARP routes).
    \item Each \textbf{cuckoo egg} is a new solution generated by a
          \emph{discrete L\'evy flight} from the solution of a nest.
    \item If the egg is \emph{better} than the nest it lands on, it replaces it
          (\emph{greedy competition}).
    \item With probability $p_a$, the \emph{worst nests} are abandoned and
          replaced by new solutions generated around the best nest
          (\emph{diversification}).
\end{itemize}

The difference with respect to ABC is that Cuckoo Search generates its candidate
solutions through the L\'evy flight, which combines many short steps with
occasional long jumps, instead of the unit perturbation of the employed bees or
the roulette-wheel selection of the onlookers.

% ------------------------------------------------------------
\section{Notation}
% ------------------------------------------------------------

\begin{itemize}
    \item $G = (V, E, R)$: graph of the CARP instance. $R \subseteq E$ is the
          set of required edges.
    \item $N = \{s_1, \ldots, s_{K}\}$: set of nests (active solutions),
          with $K = \mathit{num\_nests}$.
    \item $s^{*}$: best solution found over the whole run.
    \item $c(s)$: pure cost of solution $s$.
    \item $\mathrm{viol}(s)$: total demand excess with respect to the vehicle
          capacity.
    \item $f(s) = c(s) + \lambda \cdot \mathrm{viol}(s)$: penalized objective
          ($\lambda > 0$).
    \item $p_a \in (0,1)$: fraction of worst nests abandoned per iteration.
    \item $\beta > 0$: shape parameter of the L\'evy flight ($\beta = 1.5$
          classical).
    \item $L$: number of steps of the discrete L\'evy flight (random variable
          bounded to $[1, 12]$, see Section~\ref{sec:levy}).
    \item $P_{\mathrm{base}}$: base scale of the L\'evy steps.
    \item $\mathcal{O}_{\mathrm{intra}}$, $\mathcal{O}_{\mathrm{inter}}$:
          sets of intra-route and inter-route operators (9 operators in total).
    \item $p_{\mathrm{inter}} \in [0,1]$: probability of choosing the
          inter-route group when the solution is feasible.
    \item $p_{\mathrm{eff}}$: effective probability of choosing inter-route
          after the dynamic floor under capacity violation.
    \item $n = |R|$: number of required tasks (\emph{instance-aware} scaling
          variable).
\end{itemize}

\medskip
\noindent
\textbf{Instance-aware calibration (default values when not provided):}
\[
\mathit{iterations} = \max(200,\; 20 \cdot n), \quad
K = \max(10,\; \lceil 2\sqrt{n} \rceil), \quad
P_{\mathrm{base}} = \max\!\left(3,\; \left\lfloor\tfrac{\sqrt{n}}{2}\right\rfloor\right), \quad
M = \max(50,\; 3 \cdot n).
\]

\medskip
\noindent
\textbf{Dynamic inter/intra bias (shared with SA, simple TS, RTS and simple
ABC).} The mean violation
$\bar{v} = \tfrac{1}{K}\sum_{i=1}^{K}\mathrm{viol}(s_i)$ is computed and:
\[
p_{\mathrm{eff}} \;=\;
\begin{cases}
  \max(p_{\mathrm{inter}},\; 0.8) & \text{if } \bar{v} > 0, \\
  p_{\mathrm{inter}}              & \text{if } \bar{v} = 0.
\end{cases}
\]

% ------------------------------------------------------------
\section{Discrete L\'evy flight}
\label{sec:levy}
% ------------------------------------------------------------

In the original continuous space, the L\'evy flight produces steps whose length
follows a heavy-tailed distribution. In the discrete space of CARP routes it is
approximated by the \emph{number of consecutive perturbations} $L$:

\[
x \sim |\mathcal{N}(0,\,1)|, \qquad
L = \min\!\left(12,\; 1 + \left\lfloor x^{1/\beta} \cdot P_{\mathrm{base}} \right\rfloor\right).
\]

A large $x$ (tail of the normal distribution) produces a large $L$ and hence a
long jump in the solution space. The bound $L \leq 12$ prevents excessively
long flights that would degrade the quality of the exploration.

Each of the $L$ steps is a local perturbation: apply an operator from
$\mathcal{O}_{\mathrm{group}}$ (chosen by the inter/intra bias) to the current
solution and take the resulting solution. The concrete operator within the
group is selected uniformly at random.

% ------------------------------------------------------------
\section{Pseudocode}
% ------------------------------------------------------------

% -------- Discrete Lévy flight (auxiliary) -------------------
\begin{algorithm}[H]
\DontPrintSemicolon
\caption{DiscreteL\'evyFlight($s_{\mathrm{base}},\, P_{\mathrm{base}},\, \beta,\,
         \mathcal{O}_{\mathrm{group}}$)}
\label{alg:levy}

\KwIn{Base solution $s_{\mathrm{base}}$;\; scale $P_{\mathrm{base}}$;\;
      parameter $\beta$;\; operator group $\mathcal{O}_{\mathrm{group}}$}
\KwOut{Cuckoo solution $s'$;\; list of applied moves $\mathbf{m}$}

\BlankLine
$x \leftarrow |{\cal N}(0,1)|$
    \tcp{absolute value of a standard normal sample}
$L \leftarrow \min\!\left(12,\; 1 + \lfloor x^{1/\beta} \cdot
    \max(1, P_{\mathrm{base}}) \rfloor\right)$
    \tcp{number of flight steps; bounded to $[1,12]$}
$s' \leftarrow$ copy of $s_{\mathrm{base}}$ \;
$\mathbf{m} \leftarrow [\,]$ \;
\For{$\ell \leftarrow 1$ \KwTo $L$}{
    $(s',\, m_\ell) \leftarrow \textsc{GenerateNeighbor}(s',\, \mathcal{O}_{\mathrm{group}})$
        \tcp{local perturbation (uniform operator within the group)}
    $\mathbf{m} \leftarrow \mathbf{m} + [m_\ell]$ \;
}
\Return $(s',\, \mathbf{m})$
\end{algorithm}

\bigskip

% -------- Full Cuckoo Search ---------------------------------
\begin{algorithm}[H]
\DontPrintSemicolon
\caption{Discrete Cuckoo Search for the CARP (Yang \& Deb 2009, adapted)}
\label{alg:cuckoo}

\KwIn{CARP instance $G$;\; initial solution $s_0$;\;
      $p_a$;\; $\beta$;\; $p_{\mathrm{inter}}$;\; $\lambda$;\;
      optionally: $K$, $P_{\mathrm{base}}$, $\mathit{iterations}$,
      $M$ (all instance-aware when omitted)}
\KwOut{Best solution found $s^{*}$}

\BlankLine
\tcp{Instance-aware calibration}
$n \leftarrow |R|$ \;
$\mathit{iterations} \leftarrow \max(200,\; 20 \cdot n)$ \;
$K            \leftarrow \max(10,\; \lceil 2\sqrt{n} \rceil)$ \;
$P_{\mathrm{base}} \leftarrow \max(3,\; \lfloor\sqrt{n}/2\rfloor)$ \;
$M            \leftarrow \max(50,\; 3 \cdot n)$
    \tcp{max.\ iterations without improvement}

\BlankLine
\tcp{One-time precomputation of the intra/inter partitions}
Build $\mathcal{O}_{\mathrm{intra}}$, $\mathcal{O}_{\mathrm{inter}}$,
$\mathcal{O}_{\mathrm{fallback}}$ from the active operators \;

\BlankLine
\tcp{Nest initialization}
$N[0] \leftarrow s_0$ \;
\For{$i \leftarrow 1$ \KwTo $K - 1$}{
    $N[i] \leftarrow$ random neighbor of $N[0]$ \;
}
$s^{*} \leftarrow \arg\min_{i} f(N[i])$ \;
$\mathit{iter\_no\_improve} \leftarrow 0$ \;

\BlankLine
\For{$\mathit{cycle} \leftarrow 1$ \KwTo $\mathit{iterations}$}{

    \tcp{Early-stopping criterion by stagnation}
    \If{$\mathit{iter\_no\_improve} \geq M$}{
        \textbf{break} \;
    }
    $f^{*}_{\mathrm{start}} \leftarrow f(s^{*})$ \tcp{snapshot to detect improvement}

    \BlankLine
    \tcp{Mean violation and $p_{\mathrm{eff}}$ for this cycle}
    $\bar{v} \leftarrow \tfrac{1}{K}\sum_{i=1}^{K} \mathrm{viol}(N[i])$ \;
    \lIf{$\bar{v} > 0$}{$p_{\mathrm{eff}} \leftarrow \max(p_{\mathrm{inter}},\; 0.8)$}
    \lElse{$p_{\mathrm{eff}} \leftarrow p_{\mathrm{inter}}$}

    \BlankLine
    \tcp{\textbf{Step 1 --- Cuckoo generation through L\'evy flights}}
    \For{$i \leftarrow 0$ \KwTo $K - 1$}{
        $\mathcal{O}_{\mathrm{group}} \leftarrow$ \textsc{ChooseGroup}($p_{\mathrm{eff}},\,
            \mathrm{viol}(N[i]),\,
            \mathcal{O}_{\mathrm{inter}},\, \mathcal{O}_{\mathrm{intra}}$) \;
        $(\mathit{ck}[i],\, \mathbf{m}[i]) \leftarrow$ \textsc{DiscreteL\'evyFlight}($N[i],\,
            P_{\mathrm{base}},\, \beta,\, \mathcal{O}_{\mathrm{group}}$) \;
        \ForEach{$m \in \mathbf{m}[i]$}{\textsc{Propose}($m$)} \;
    }

    \BlankLine
    \tcp{Vectorized batch evaluation of the cuckoos (GPU when available)}
    $\{f(\mathit{ck}[i])\}_{i=0}^{K-1} \leftarrow$ \textsc{EvaluateBatch}($\{\mathit{ck}[i]\}$) \;

    \BlankLine
    \tcp{\textbf{Step 2 --- Competition: cuckoo vs.\ random nest}}
    \For{$i \leftarrow 0$ \KwTo $K - 1$}{
        $j \leftarrow \textsc{UniformInteger}(0,\, K-1)$
            \tcp{rival nest chosen at random}
        \If{$f(\mathit{ck}[i]) < f(N[j]) - \varepsilon$}{
            $N[j] \leftarrow \mathit{ck}[i]$
                \tcp{the cuckoo replaces the rival nest}
            \textsc{Accept}($\mathbf{m}[i][-1]$) \;
            \tcp{Immediate detection of global improvement (imputation-bug fix)}
            \If{$f(\mathit{ck}[i]) < f(s^{*}) - \varepsilon$}{
                $s^{*} \leftarrow \mathit{ck}[i]$ \;
                \textsc{RecordImprovement}($\mathbf{m}[i][-1]$)
                    \tcp{operator of the last L\'evy step of THIS cuckoo}
            }
        }
    }

    \BlankLine
    \tcp{\textbf{Step 3 --- Abandonment of the worst nests}}
    $n_a \leftarrow \max(1,\; \lfloor p_a \cdot K \rfloor)$
        \tcp{number of nests to abandon}
    $\mathit{worst} \leftarrow$ indices of the $n_a$ nests with largest $f(\cdot)$ \;
    $j^{*} \leftarrow \arg\min_{i} f(N[i])$
        \tcp{index of the best nest}
    $\mathcal{O}_{\mathrm{ab}} \leftarrow$ \textsc{ChooseGroup}($p_{\mathrm{eff}},\,
        \mathrm{viol}(N[j^{*}]),\,
        \mathcal{O}_{\mathrm{inter}},\, \mathcal{O}_{\mathrm{intra}}$) \;
    \ForEach{$k \in \mathit{worst}$}{
        $(N[k],\, \mathbf{m}_{\mathrm{ab}}) \leftarrow$ \textsc{DiscreteL\'evyFlight}($N[j^{*}],\,
            P_{\mathrm{base}},\, \beta,\, \mathcal{O}_{\mathrm{ab}}$)
            \tcp{L\'evy flight from the current best nest}
        \ForEach{$m \in \mathbf{m}_{\mathrm{ab}}$}{\textsc{Propose}($m$)} \;
        \tcp{(evaluation included in the final block)}
    }

    \BlankLine
    \tcp{Vectorized evaluation of the freshly replaced nests}
    Evaluate $f(N[k])$ for all $k \in \mathit{worst}$ \;
    \ForEach{$k \in \mathit{worst}$}{
        \If{$f(N[k]) < f(s^{*}) - \varepsilon$}{
            $s^{*} \leftarrow N[k]$ \tcp{abandonment may also improve the global best}
        }
    }

    \BlankLine
    \tcp{Stagnation-counter update}
    \lIf{$f(s^{*}) < f^{*}_{\mathrm{start}} - \varepsilon$}{%
        $\mathit{iter\_no\_improve} \leftarrow 0$}
    \lElse{%
        $\mathit{iter\_no\_improve} \leftarrow \mathit{iter\_no\_improve} + 1$}
}

\BlankLine
\Return $s^{*}$
\end{algorithm}

% ------------------------------------------------------------
\section{Implementation notes}
% ------------------------------------------------------------

\paragraph{Vectorized evaluation (GPU).}
In Step 1, the $K$ cuckoos are evaluated in a single batch with
\texttt{costo\_lote\_penalizado\_ids} (NumPy on CPU; CuPy on GPU when
\texttt{usar\_gpu=True}). The nests replaced in Step 3 are also evaluated in a
batch. This concentrates the vectorizable work into two calls per iteration,
minimizing pure-Python overhead.

\paragraph{Imputation-bug fix.}
In the previous version, the detection of a global improvement was performed by
comparing the best value before and after the whole iteration, crediting the
\emph{last accepted move} of the iteration (which might not be the responsible
operator). In this version, the call \textsc{RecordImprovement}$(m)$ happens
inside the same \texttt{if} that detects $f(\mathit{ck}[i]) < f(s^{*})$, with
the operator of the last L\'evy step of \emph{that} specific cuckoo.
Consequence: $\sum_{\mathrm{op}}\mathit{improved}[\mathrm{op}] = \mathit{improvements}$.

\paragraph{Nest initialization.}
Nest 0 is seeded with the best available initial solution. Nests $1 \ldots K-1$
are generated as random neighbors of nest 0 (same strategy as in simple ABC),
which provides immediate diversity without discarding the information of the
starting solution. In the val/egl campaign, the initial solution is produced by
a Path-Scanning warm start (best of 5 constructions).

\paragraph{Difference with simple ABC in the scouts.}
In Cuckoo Search, the abandonment step (Step 3) replaces the worst nests with a
L\'evy flight \emph{from the current best nest} (guided exploration). In simple
ABC, scouts replace exhausted sources with completely random solutions (pure
exploration). Cuckoo Search tends to explore around the best-known region;
simple ABC keeps a higher global diversification.

\paragraph{Parameter $\beta$.}
The classical value of Yang \& Deb is $\beta = 1.5$. Smaller values (e.g.,
$\beta = 1.2$) yield heavier-tailed distributions: more long jumps, more
exploration but less local precision. Larger values (e.g., $\beta = 2.0$)
produce shorter steps, favoring exploitation. In MetaCARP, the calibrated value
$\beta = 1.3$ is used as default for the experimental campaigns.

\end{document}
